\documentclass[11pt,a4paper,twoside,openright]{book}
\usepackage[ngerman]{babel}
\newcommand{\HandbookPdfTitle}{SASD Math Toolkit - CSharp/.NET Handbuch Version 1.0}
\newcommand{\HandbookFooter}{Handbuch Version 1.0 - Deutsch}
\newcommand{\HandbookSubtitle}{Numerical Methods - C\# / .NET}
\newcommand{\HandbookVersionText}{Benutzer- und Entwicklerhandbuch - Version 1.0}
\newcommand{\HandbookTagline}{Numerische Verfahren für robuste, nachvollziehbare und wiederverwendbare C\#-Anwendungen.}
\newcommand{\HandbookRepositoryRef}{Repository: \repo\quad |\quad C\#/.NET 10}

\usepackage{sasd-handbook}

\begin{document}
\frontmatter
\SASDTitlePage

\chapter*{Dokumentstatus}
\addcontentsline{toc}{chapter}{Dokumentstatus}
\begin{sasdnote}{Version 1.0 dieses Handbuchs}
Dieses Dokument ist die deutschsprachige Handbuchausgabe 1.0 für die C\#/.NET-Implementierung des SASD Math Toolkit. Die Kapitelinhalte werden beim Build aus dem gepflegten Markdown-Handbuch unter \texttt{docs/user-guide/de/} erzeugt. Dadurch bleibt Markdown die redaktionelle Quelle, während LaTeX Satz, Navigation und PDF-Ausgabe übernimmt.
\end{sasdnote}

\begin{tabularx}{\textwidth}{@{}>{\bfseries}p{0.31\textwidth}X@{}}
Dokument & SASD Math Toolkit - C\#/.NET Benutzer- und Entwicklerhandbuch\\
Handbuchversion & 1.0\\
Repository & \repo\\
Zielplattform & .NET 10\\
Assembly & \texttt{Sasd.Math.Toolkit}\\
Root Namespace & \texttt{Sasd.Numerics}\\
Lizenz & MIT\\
Autor / Projekt & Robin Goerlach; SASD-GmbH\\
\end{tabularx}

\chapter*{Vorwort}
\addcontentsline{toc}{chapter}{Vorwort}
Das SASD Math Toolkit überträgt den funktionalen Kern einer klassischen Numerical-Methods-Toolbox in eine moderne, unabhängig entwickelte C\#/.NET-Bibliothek. Das historische Handbuch dient dabei als Orientierung für Problemklassen und Handbuchdramaturgie - nicht als Text- oder Quellcodevorlage.

Die historische Ausgabe strukturierte numerische Verfahren konsequent nach Problem, Beschreibung, Eingaben, Ausgaben, Aufrufsyntax, Kommentaren und Beispiel. Die vorliegende Ausgabe übernimmt diesen anwenderorientierten Gedanken, formuliert ihn aber für moderne C\#-APIs mit Exceptions, Statusobjekten, Residuen und Tests neu.

\begin{sasdtip}{Leitgedanke}
Eine Zahl ist noch kein belastbares numerisches Ergebnis. Für wichtige Rechnungen gehören Methode, Toleranz, Terminierungsstatus und eine problemspezifische Qualitätsprüfung zusammen.
\end{sasdtip}

\tableofcontents
\clearpage
\mainmatter

\input{build/de/erste-schritte.tex}
\input{build/de/nullstellen.tex}
\input{build/de/interpolation.tex}
\input{build/de/differentiation.tex}
\input{build/de/integration.tex}
\input{build/de/matrizen-lineare-gleichungssysteme.tex}
\input{build/de/eigenwerte-eigenvektoren.tex}
\input{build/de/differentialgleichungen.tex}
\input{build/de/least-squares.tex}
\input{build/de/fft-faltung-korrelation.tex}
\input{build/de/diagnostik-fehlerfaelle.tex}

\appendix
\chapter{Historischer Kontext und V1-Abdeckung}
Die Kapitelreihenfolge der C\#-Ausgabe folgt bewusst den großen numerischen Problemklassen der historischen Numerical Methods Toolbox. Die Umsetzung ist jedoch keine Quellcode-Portierung. Die öffentlichen APIs, Datentypen, Fehlerbehandlung, Beispiele und Diagnosen gehören zum SASD Math Toolkit.

\begin{longtable}{@{}p{0.28\textwidth}p{0.62\textwidth}@{}}
\toprule
\textbf{Historischer Themenblock} & \textbf{C\#/.NET-Ausgabe}\\
\midrule
Nullstellen & Bisektion, Newton-Raphson, Sekante, Newton-Horner, Muller, Laguerre und Polynomdeflation\\
Interpolation & Lagrange, dividierte Differenzen, natürliche und geklemmte kubische Splines\\
Differentiation / Integration & Funktions-, Tabellen- und Spline-Differentiation; Newton-Cotes, adaptiv, Romberg, Gauß-Legendre\\
Matrixverfahren & \texttt{DenseMatrix}, Gauß, pivotierte LU, Householder-QR, Inverse, Residuen und Gauss-Seidel\\
Eigenwertprobleme & Potenz-/inverse Potenzmethode, Wielandt-Deflation und zyklisches Jacobi-Verfahren\\
Differentialgleichungen & RK4-Familie, RKF45, Adams-Prädiktor-Korrektor sowie lineares und nichtlineares Shooting\\
Least Squares & Polynom- und Basisfits, benannte Modelle; modernisiert auf Householder-QR\\
FFT / Anwendungen & Komplexe und reelle FFT, kompaktes Realspektrum, Faltung und Kreuzkorrelation\\
Grafische Demonstrationen & Plattformunabhängiges, deterministisches HTML/SVG-Sample statt Kopplung der Numerik an eine GUI-Technologie\\
\bottomrule
\end{longtable}

\chapter{Build und Pflege der LaTeX-Ausgabe}
Die LaTeX-Ausgabe wird absichtlich aus den bestehenden Markdown-Kapiteln generiert. Änderungen an fachlichen Erläuterungen sollen daher zuerst in \texttt{docs/user-guide/de/} bzw. \texttt{docs/user-guide/en/} erfolgen. Das Verzeichnis \texttt{LaTeX/} enthält Layout, Buildlogik und sprachspezifische Front-/Back-Matter.

\begin{lstlisting}
cd docs/user-guide/LaTeX
make de      # deutsches PDF
make en      # English PDF
make all     # beide Ausgaben
\end{lstlisting}

\backmatter
\chapter*{Lizenz und Herkunft}
\addcontentsline{toc}{chapter}{Lizenz und Herkunft}
Das SASD Math Toolkit steht unter der MIT-Lizenz. Das historische Borland-Handbuch und die historischen Produktnamen bleiben Eigentum ihrer jeweiligen Rechteinhaber. Dieses Handbuch dokumentiert die unabhängig entwickelte SASD-C\#/.NET-Implementierung und reproduziert nicht den historischen Handbuchtext.

\end{document}
